\section*{Notation and conventions}
\addcontentsline{toc}{section}{Notation and conventions}
\label{sec:notation}

Every symbol used in this report is listed here. Each is defined once, in the
section given in the last column, and is used with that meaning everywhere else.

\medskip
\noindent\begin{tabular}{@{}l>{\raggedright\arraybackslash}p{8.9cm}l@{}}
\toprule
Symbol & Meaning & Defined in \\
\midrule
$n$              & number of jobs                                              & \S\ref{ssec:ssp} \\
$J$              & the set of jobs, $J=\rset{1,\dots,n}$                        & \S\ref{ssec:ssp} \\
$T$              & the set of tools                                            & \S\ref{ssec:ssp} \\
$\lvert T\rvert$ & number of tools                                             & \S\ref{ssec:ssp} \\
$T_j$            & the tools job $j$ requires                                   & \S\ref{ssec:ssp} \\
$b$              & magazine capacity, in tool slots                            & \S\ref{ssec:ssp} \\
$\U$             & the tools required by at least one job, $\U=\bigcup_j T_j$   & \S\ref{ssec:ssp} \\
  $q$              & the tool excess, $q=\lvert\U\rvert-b\ge0$ after preprocessing & \S\ref{ssec:ssp} \\
$\pi$            & a job order (a permutation of $J$)                           & \S\ref{ssec:ssp} \\
$M_i$            & the magazine contents at position $i$                        & \S\ref{ssec:ssp} \\
  $\Zst$           & optimal free-initial switch cost                              & \S\ref{ssec:ssp} \\
  $\Zempty$        & optimal empty-start switch cost                               & \S\ref{ssec:ssp} \\
$\KTNS(\pi)$     & the cheapest loading cost for the fixed order $\pi$          & \S\ref{ssec:ktns} \\
$G$              & a group of jobs                                             & \S\ref{ssec:jgp} \\
$\Kst$           & the fewest feasible groups covering all jobs                 & \S\ref{ssec:jgp} \\
  $C$              & a configuration: a $b$-tool subset of $\U$                  & \S\ref{ssec:gtsp} \\
$\confs$         & the set of all configurations                                & \S\ref{ssec:gtsp} \\
$H_j$            & the cluster of configurations able to serve job $j$          & \S\ref{ssec:gtsp} \\
$d(C,C')$        & transition cost, $d(C,C')=\lvert C'\setminus C\rvert$        & \S\ref{ssec:gtsp} \\
$\mathcal{P}$    & a grouping: a partition of $J$ into feasible groups          & \S\ref{ssec:frames} \\
$\gamma(\mathcal{P})$ & the cheapest walk through the configurations of $\mathcal{P}$ & \S\ref{ssec:exact} \\
$\Hwalk$         & the walk value of the two-phase heuristic                    & \S\ref{ssec:frames} \\
$H$              & the value the two-phase heuristic returns                    & \S\ref{ssec:frames} \\
$R$              & re-insertions: insertions of tools evicted earlier           & \S\ref{ssec:ratio} \\
$R_{\mathrm{opt}}$ & re-insertions made by an optimal schedule                  & \S\ref{ssec:ratio} \\
$w_{ij}$         & pairwise bound, $w_{ij}=\max(0,\lvert T_i\cup T_j\rvert-b)$  & \S\ref{ssec:bbc} \\
$\rho$           & setup cost charged per configuration change                  & \S\ref{ssec:collapse} \\
$\theta$         & the Benders stand-in for the loading cost                    & \S\ref{ssec:bbc} \\
\bottomrule
\end{tabular}

\subsection*{Two conventions that are fixed once and never varied}

\paragraph{Cost.}
Formulations in the literature charge the first magazine fill differently, and
comparing their objective values without adjusting for this produces differences
that are not real. Two costs are used here and are always named.
The \emph{empty-start} cost counts every insertion, beginning from an empty
magazine. The \emph{free-initial} cost does not charge the first fill.
Proposition~\ref{prop:conv} shows the two differ by a constant that depends only on
the instance. Section~\ref{ssec:conv} gives the details. Theory is stated in the
free-initial cost; the computational study reports the empty-start cost, because
that is what the solvers return.

\paragraph{Status of results.}
Every formal result and every quantitative claim carries a label saying what kind of
statement it is.

\begin{itemize}[topsep=3pt,itemsep=1pt]
  \item \textbf{Theorem}, \textbf{Proposition}, \textbf{Lemma}, \textbf{Corollary}
  --- proved, with the proof given here or in Appendix~\ref{app:proofs}.
  \item \textbf{Observation} --- established by computation, not by proof. The
  computation is named, and Appendix~\ref{app:verify} says how to repeat it.
  \item \textbf{Conjecture} --- supported by computation but not proved. Where a
  partial argument exists it is given under the heading \emph{Supporting
  argument}, which is deliberately not the word \emph{Proof}.
\end{itemize}

\noindent
Numbers taken from the literature name the paper and the table or section they
come from. Numbers computed for this report name the instance family, its
parameter range, and the script that produced them.
